\documentclass{article}
\usepackage{graphicx, enumitem, parskip} % Required for inserting images
\setlist[itemize]{label=-}


\title{Casella, Berger - Chapter 1, Exercise 1.2}
\author{Afonso Vilhena Francisco}
\date{24 November 2026}

\begin{document}

\maketitle

$\mathbf{1.2}$ Verify the following identities.
\begin{enumerate}[label=(\alph*)]
    \item $A\backslash B = A \backslash (A \cap B) = A \cap B^c$
    \item $B = (B \cap A) \cup (B \cap A^c)$
    \item $B\backslash A = B \cap A^c$
    \item $A \cup B = A \cup (B \cap A^c)$
\end{enumerate}

\vspace{1cm}
\begin{enumerate}[label=(\alph*)]
    \item Let $x \in A\backslash B$, that is, $x \in A$ and $x \notin B$. Since $x \notin B$, then $x \notin A \cap B$. Since $x \in A$ and $x \notin A \cap B$, then $x \in A \backslash (A \cap B)$.\\
    Conversely, let $x \in A \backslash (A \cap B)$, then $x \in A$ and $x \notin (A \cap B)$. The latter means that $x \notin A$ or $x \notin B$. Since, $x \in A$, then $x \notin B$. Since, $x \in A$ and $x \notin B$, then $x \in A \backslash B$.\\
    
    \par
    
    Let $x \in A \backslash B$, then $x \in A$ and $x \notin B$. Since $x \notin B$, then $x \in B^c$. Since $x \in A$ and $x \in B^c$, then $x \in A \cap B^c$.\\
    Conversely, let $x \in A\cap B^c$, then $x \in A$ and $x \in B^c$. Since $x \in B^c$, then $x \notin B$. Since $x \in A$ and $x \notin B$, then $x \in A \backslash B.$

    \vspace{0.5cm}
    
    \item Let $x$ in $B$. 
    \begin{itemize}
        \item $(x \in A)$ then $x \in B \cap A$ and then $x \in (B \cap A) \cup (B \cap A^c)$.
        \item $(x \notin A)$ then $x \in A^c$, then $x \in B \cap A^c$, then $x \in (B \cap A) \cup (B \cap A^c)$
    \end{itemize}

    Conversly, let $x \in (B \cap A) \cup (B \cap A^c),$ then $x \in (B \cap A)$ or $x \in (B \cap A^c)$.
    \begin{itemize}
        \item $(x \in (B \cap A))$ then $x \in B$ and $x \in A$, therefore $x \in B$.
        \item ($x \in (B \cap A^c))$, then $x \in B$ and $x \in A^c$, therefore $x \in B$.
    \end{itemize}

    \vspace{0.5cm}

    \item Interchanging the roles of $A$ and $B$ in $(a)$, we obtain the pretended result.

    \vspace{0.5cm}

    \item Let $x \in A \cup B$, then $x \in A$ or $x \in B$.
    \begin{itemize}
        \item ($x \in A$) then $x \in A$ or $x \in (B \cap A^c),$ therefore $x \in A \cup(B \cap A^c)$.
        \item ($x \in B$)
        \begin{itemize}
            \item ($x \in A$) then $x \in A$ or $x \in (B \cap A^c)$, therefore $x \in A \cup (B \cap A^c)$.
            \item ($x \notin A$) then $x \in B$ and $x \in A^c$, then $x \in (B \cap A^c)$, then $x \in A$ or $x \in (B \cap A^c)$, therefore $x \in A \cup (B \cap A^c)$.
        \end{itemize}
    \end{itemize}

    Conversly, let $x \in A \cup (B \cap A^c)$, then $x \in A$ or $x \in (B \cap A^c)$.
    \begin{itemize}
        \item ($x \in A$) then $x \in A$ or $x \in B$, therefore $x \in (A \cup B)$.
        \item ($x \in (B \cap A^c)$) then $x \in B$ and $x \in A^c$, then $x \in B$, then $x \in A$ or $x \in B$, therefore $x \in (A \cup B)$.
    \end{itemize}
\end{enumerate}



\end{document}
